 \begin{lemma}
 \label{middle 5}
 Suppose $G$ contains a subgraph of the form shown in Figure
~\ref{mid5 graph}. Then $W$ is $\la(2)$-infinite. 
 
 \begin{figure}[htb]\centering
 \begin{tikzpicture}
 \node[main node] (1) {};
 \node[main node] (2) [right=1cm of 1] {};
 \node[main node] (3) [right=1cm of 2] {};
 \node[main node] (6) [right=1cm of 3] {};

 \node (11) [below=0.2cm of 1] {$a$};
 \node (22) [below=0.1cm of 2] {$b$};
 \node (33) [below=0.2cm of 3] {$c$};
 \node (66) [below=0.1cm of 6] {$d$};

 \path[draw]
 (1) edge node {} (2)
 (2) edge node [above] {$5$} (3)
 (3) edge node {} (6);
 \end{tikzpicture}
 \caption{}
 \label{mid5 graph}
 \end{figure}
 \end{lemma}

 \begin{proof}
 Let $X=acbc$ and $Y=bdcb$. For each $k\in \Z_{\ge 0}$, let $w_k=XY\cdots$ be
 the string that starts in
 $X$ and contains
 $k$ alternating occurrences of $X$ and $Y$. The heap of $w_k$ is shown
 in Figure~\ref{mid5 heap}, where the dashed rectangles alternately correspond
 to the
 expression $X$ and $Y$ and appear a total of $k$ times. Observe that
 $w_k$ is reduced for each $k$ by Proposition~\ref{fc heap criterion}.
 \begin{figure}[htb]\centering
 \resizebox{0.4\textwidth}{!}{
 \begin{tikzpicture}
 %
 %
 \node[main node] (3) {};
 \node[main node] (4) [right = 3cm of 3] {};
 \node[main node] (5) [above right=0.75cm and 1.5cm of 3] {};
 \node[main node] (6) [above right=0.75cm and 1.5cm of 5] {};
 \node[main node] (7) [above left=0.75cm and 1.5cm of 6] {};
 \node[main node] (8) [above right=0.75cm and 1.5cm of 6] {};
 \node[main node] (9) [above right=0.75cm and 1.5cm of 7] {};
 \node[main node] (10) [above left=0.75cm and 1.5cm of 9] {};
 \node[main node] (11) [above left=0.75cm and 1.5cm of 10] {};
 \node[main node] (12) [above right=0.75cm and 1.5cm of 10] {};
 \node[main node] (13) [above right=0.75cm and 1.5cm of 11] {};
 \node[main node] (14) [above right=0.75cm and 1.5cm of 13] {};
 \node[main node] (15) [above left=0.75cm and 1.5cm of 14] {};
 \node[main node] (16) [above right=0.75cm and 1.5cm of 14] {};
 \node (17) [above right=0.75cm and 1.5cm of 15] {};
 \node (18) [above left=0.75cm and 1.5cm of 17] {};

 %
 %
 \node (33) [above=0.01cm of 3] {$a$};
 \node (44) [above=0.01cm of 4] {$c$};
 \node (55) [above=0.01cm of 5] {$b$};
 \node (66) [above=0.01cm of 6] {$c$};
 \node (77) [above=0.01cm of 7] {$b$};
 \node (88) [above=0.01cm of 8] {$d$};
 \node (99) [above=0.01cm of 9] {$c$};
 \node (1010) [above=0.01cm of 10] {$b$};
 \node (1111) [above=0.01cm of 11] {$a$};
 \node (1212) [above=0.01cm of 12] {$c$};
 \node (1313) [above=0.01cm of 13] {$b$};
 \node (1414) [above=0.01cm of 14] {$c$};
 \node (1515) [above=0.01cm of 15] {$b$};
 \node (1616) [above=0.01cm of 16] {$d$};
 %

 \path[draw]
 %
 (3)--(5)--(6)--(8)--(9)--(10)--(11)--(13)--(14)--(16) %
 (4)--(5)
 (6)--(7)--(9)
 (10)--(12)--(13)
 (14)--(15);

 \path[draw,dashed]
 (15)--(17)--(16);

 \node (r1) [below left=0.2cm and 0.5cm of 3] {};
 \node (r2) [below right=0.2cm and 2cm of 4] {};
 \node (r3) [above = 2.4cm of r1] {};
 \node (r4) [above = 2.4cm of r2] {};
 \node (r5) [above = 2.3cm of r3] {};
 \node (r6) [above = 2.3cm of r4] {};
 \node (r7) [above = 2.2cm of r5] {};
 \node (r8) [above = 2.2cm of r6] {};
 \node (r9) [above = 2.3cm of r7] {};
 \node (r10) [above = 2.3cm of r8] {};
 \node (r11) [above right=1.3cm and 0.5cm of 15] {$\vdots$};


 \path[draw,dashed,blue]
 (r1)--(r2)--(r4)--(r3)--(r1)
 (r4)--(r6)--(r5)--(r3)
 (r5)--(r7)--(r8)--(r6)
 (r7)--(r9)
 (r8)--(r10);

 \end{tikzpicture}
 }
 %
 %
 %
 \caption{}
 \label{mid5 heap}
 \end{figure}

 We shall prove that
 \begin{equation}
 \label{eq:even}
 w_kac\le_R w_{k+2}c\le_R w_{k+2}\le_R w_{k+1} \le_R w_kac
 \end{equation}
 for all even integers $k\ge 0$ and that
 \begin{equation}
 \label{eq:odd}
 w_kbd\le_R w_{k+2}b\le_R w_{k+2}\le_R w_{k+1}\le_R w_kbd
 \end{equation}
 for all odd integers $k\ge 1$. It then follows that $w_k\sim_R ac$ and hence
 $\la(w_k)=\la(ac)=2$ for all $k\ge 1$, therefore $W$ is $\la(2)$-infinite. 
 



 To prove~\eqref{eq:even}, let $k\ge 0$ be an even integer. Note that
 $w_{k+2}c\le_R w_{k+2}\le_R w_{k+1} \le_R w_kac$ follows from Corollary
~\ref{weak order}, therefore it suffices to show that $w_kac\le_R
 w_{k+2}c$. Let $x=w_kac$ and $y=w_{k+2}c$. We will show in fact
 $x\prec_R y$. Since $x<y$ and $a\in
 \mathcal{R}(x)\setminus\mathcal{R}(y)$, it
 further suffices to show that $\mu(x,y)\neq 0$ by Proposition~\ref{alternative prec}.

 To compute $\mu(x,y)$, we consider the coset decompositions of $x$ and
 $y$ with respect to $I=\{b,c\}$, where 
\begin{IEEEeqnarray*} {RRLLLLLLLL}
 x&=& \cdots bdcbac = x^I\cdot {x}_I  &\quad & \text{with}\quad 
 x^I&= w_ka,
 & x_I&= c,\\
 y&=& \cdots acbcdbcbc= y^I\cdot y_I  &\quad & \text{with} \quad
 y^I&=w_{k+1}d, &{}^I
 y&= bcbc.
 \end{IEEEeqnarray*}
 For any integers $0\le i,j \le 4$, let $p_i$ be the word $cb\cdots$ that
 alternates in $b$ and $c$, starts in $c$, and has length $i$, and similarly let
 $q_j$ be the alternating word $bc\cdots$ of length $j$. Let
 $x_i=x^I\cdot p_i$ and $y_j=y^I\cdot q_j$, and set
 \[
 [i,j]=\mu(x_i, y_j)
 \]
 for all $0\le i,j\le 4$. We have
 \begin{align*}
 [1,4]&= -[3,4]+[2,3]\\
 &= -[3,4]+(-[4,3]+[3,2]+[3,4])\\
 &= -[4,3]+[3,2]\\
 &= -[4,3]+([4,1]+[4,3])\\
 &= [4,1],
 \end{align*}
 where the first, second and fourth equalities follow from applications of Part~\eqref{prop3.7_1} of
 Proposition~\ref{star relation} with $(x_2,y_4), (x_3,y_3)$ and $(x_4,y_2)$ in
 place of the pair $(x,y)$, respectively. Now, 
 \[
 [4,1]=\mu(x_4,y_1)=\mu(w_kacbcb,w_kacbcdb)=1,
 \]
\looseness-1
 where the last equality follows from Corollary~\ref{difference 1},
 therefore $\mu(x,y)=[1,4]=[4,1]\neq 0$, and we have proved~\eqref{eq:even}. 

 The proof of~\eqref{eq:odd} is similar to that of $\eqref{eq:even}$,
 thanks to the symmetry $a\leftrightarrow d, b\leftrightarrow c$ in
 the subgraph in question. We have now completed our proof. 
 \end{proof}
